\chapter{Packet Sniffing and Wireshark}
\chapterguide{A basic understanding of hosts, switches, IP traffic, and why packets move through an interface.}{Define packet sniffing and a network traffic trace, navigate the Wireshark interface, capture packets, and connect a packet-list entry to protocol details and raw bytes.}
\sourcelab{Lab 1 - Wireshark.pptx}

\section{What this lab is building}
The first lab is about \emph{visibility}. Before configuring routes, services, VLANs, or WAN links, you need a way to observe what the network actually carries. The supplied slides introduce three linked ideas:
\begin{description}
  \item[Packet sniffing] tapping or observing packets while they flow across a network.
  \item[Network traffic trace] a recording of packets received and transmitted by a network interface.
  \item[Wireshark] a graphical packet analyzer used in the course to inspect packets, troubleshoot network problems, and investigate suspicious traffic.
\end{description}

\begin{conceptmap}
\textbf{Generate traffic} $\rightarrow$ \textbf{capture on the correct interface} $\rightarrow$ \textbf{find the relevant packet} $\rightarrow$ \textbf{decode protocol fields} $\rightarrow$ \textbf{relate the packet to the action that produced it}.
\end{conceptmap}

\sourceimage[0.64\linewidth]{Packet-sniffing illustration from the supplied slide}{wireshark-sniffing.png}

The illustration emphasizes the reason capture placement matters: packets must be observable on the interface or network segment being monitored. A capture taken elsewhere may be valid but irrelevant to the event you are trying to analyze.

\section{A traffic trace is evidence, not just a screenshot}
A trace is useful because it preserves a time-ordered sequence of network events. In later labs, the same idea becomes a troubleshooting method: if a ping fails, if a routing adjacency is missing, or if an application session does not establish, the trace can show what was transmitted and what response, if any, returned.

\compactsourceimage[0.64\linewidth]{Network-traffic trace concept used in the source material}{wireshark-traffic-trace.png}

\section{The Wireshark workspace}
The ``Basic UI'' slide labels the main working areas directly. The exact version of Wireshark may look different, but the analysis pattern is the same.

\sourceimage[0.94\linewidth]{Wireshark packet list, packet details, and packet bytes panes}{wireshark-ui.png}

\begin{tabularx}{\linewidth}{>{\bfseries}p{0.22\linewidth}X}
\toprule
Pane & What you use it for\\\midrule
Packet list & One row per captured packet/frame. Typical columns include time, source, destination, protocol, length, and an information summary.\\
Packet details & A decoded protocol tree for the selected packet. Expand layers to inspect addresses, control flags, protocol fields, and encapsulation.\\
Packet bytes & The raw captured bytes, normally shown in hexadecimal with an ASCII representation beside them.\\\bottomrule
\end{tabularx}

\section{A disciplined capture workflow}
\begin{netsteps}
\item Identify the traffic you want to observe and the interface that should carry it.
\item Start packet capture on that interface.
\item Generate a small, known event such as a ping so you have a clear reference point.
\item Find the corresponding packet in the packet list using addresses, protocol, and timing.
\item Expand the relevant protocol layers in Packet Details.
\item Compare decoded fields with Packet Bytes only when byte-level evidence is required.
\item Stop the capture once enough evidence has been collected; a smaller trace is easier to reason about.
\end{netsteps}

\begin{workedexample}
Suppose Host A pings Host B. In Wireshark, locate ICMP traffic and identify the request and reply. Check the source/destination addresses in each direction and confirm that the reply reverses the endpoints. The important result is not ``Wireshark shows packets''; it is that the trace proves which host sent the request, whether the peer answered, and how the protocol exchange unfolded.
\end{workedexample}

\section{What to record in a lab submission}
A useful packet-analysis screenshot should make the evidence readable. Capture enough of the interface to show the selected packet and the relevant decoded fields, and annotate or explain what the packet proves. If a filter is used, include the filter expression or describe it in the report.

\begin{verifybox}
For this lab, convincing evidence is a capture in which you can point to a known traffic event and identify at least the source, destination, protocol, and the relevant request/reply relationship.
\end{verifybox}

\begin{pitfall}
A technically successful capture can still be useless if it was taken on the wrong interface or contains too much unrelated traffic. Generate controlled traffic and narrow the view with protocol or address filters instead of searching blindly through a large capture.
\end{pitfall}

\begin{quickcheck}
\begin{enumerate}
\item What is the difference between the packet list and packet details panes?
\item Why is a traffic trace easier to interpret when you deliberately generate known traffic?
\item Which pane would you use when you need the exact bytes carried in a frame?
\item If no ping packet appears, which should you question first: the capture location or the packet bytes?
\end{enumerate}
\end{quickcheck}

\begin{chapterrecap}
Packet sniffing provides visibility, a traffic trace records the evidence, and Wireshark organizes that evidence into packet summaries, decoded protocol fields, and raw bytes. This observation skill is reused throughout the course whenever a configuration must be verified or debugged.
\end{chapterrecap}
